%\documentclass[11pt, letterpaper]{article} % for ArXiV
%\documentclass[twoside,leqno]{article}   % for SODA


\documentclass[anonymous,a4paper,USenglish,cleveref,autoref,thm-restate] {socg-lipics-v2021}
%This is a template for producing LIPIcs articles. 
%See lipics-v2021-authors-guidelines.pdf for further information.
%for A4 paper format use option "a4paper", for US-letter use option "letterpaper"
%for british hyphenation rules use option "UKenglish", for american hyphenation rules use option "USenglish"
%for section-numbered lemmas etc., use "numberwithinsect"
%for enabling cleveref support, use "cleveref"
%for enabling autoref support, use "autoref"
%for anonymousing the authors (e.g. for double-blind review), add "anonymous"
%for enabling thm-restate support, use "thm-restate"
%for enabling a two-column layout for the author/affilation part (only applicable for > 6 authors), use "authorcolumns"
%for producing a PDF according the PDF/A standard, add "pdfa"

%\pdfoutput=1 %uncomment to ensure pdflatex processing (mandatatory e.g. to submit to arXiv)
%\hideLIPIcs  %uncomment to remove references to LIPIcs series (logo, DOI, ...), e.g. when preparing a pre-final version to be uploaded to arXiv or another public repository

%\graphicspath{{./graphics/}}%helpful if your graphic files are in another directory

\bibliographystyle{plainurl}% the mandatory bibstyle
\title{Adaptive and Approximate Orthogonal Range Counting in 3-d}


%\titlerunning{Dummy short title} %TODO optional, please use if title is longer than one line

\author{Jane {Open Access}}{Dummy University Computing Laboratory, [optional: Address], Country \and My second affiliation, Country \and \url{http://www.myhomepage.edu} }{johnqpublic@dummyuni.org}{https://orcid.org/0000-0002-1825-0097}{(Optional) author-specific funding acknowledgements}%TODO mandatory, please use full name; only 1 author per \author macro; first two parameters are mandatory, other parameters can be empty. Please provide at least the name of the affiliation and the country. The full address is optional. Use additional curly braces to indicate the correct name splitting when the last name consists of multiple name parts.

\author{Joan R. Public\footnote{Optional footnote, e.g. to mark corresponding author}}{Department of Informatics, Dummy College, [optional: Address], Country}{joanrpublic@dummycollege.org}{[orcid]}{[funding]}

\authorrunning{J. Open Access and J.\,R. Public} %TODO mandatory. First: Use abbreviated first/middle names. Second (only in severe cases): Use first author plus 'et al.'

\Copyright{Jane Open Access and Joan R. Public} %TODO mandatory, please use full first names. LIPIcs license is "CC-BY";  http://creativecommons.org/licenses/by/3.0/

%\ccsdesc[100]{\textcolor{red}{Replace ccsdesc macro with valid one}} %TODO mandatory: Please choose ACM 2012 classifications from https://dl.acm.org/ccs/ccs_flat.cfm 

\ccsdesc[500]{Theory of computation~Data structures and algorithms for data management}
\ccsdesc[100]{Theory of computation~Computational geometry}
%\ccsdesc[100]{Theory of computation~Data structures design and analysis}

\keywords{Data Structures, Range Searching, Orthogonal Range Counting} %TODO mandatory; please add comma-separated list of keywords

\category{} %optional, e.g. invited paper

\relatedversion{} %optional, e.g. full version hosted on arXiv, HAL, or other respository/website
%\relatedversiondetails[linktext={opt. text shown instead of the URL}, cite=DBLP:books/mk/GrayR93]{Classification (e.g. Full Version, Extended Version, Previous Version}{URL to related version} %linktext and cite are optional

%\supplement{}%optional, e.g. related research data, source code, ... hosted on a repository like zenodo, figshare, GitHub, ...
%\supplementdetails[linktext={opt. text shown instead of the URL}, cite=DBLP:books/mk/GrayR93, subcategory={Description, Subcategory}, swhid={Software Heritage Identifier}]{General Classification (e.g. Software, Dataset, Model, ...)}{URL to related version} %linktext, cite, and subcategory are optional

%\funding{(Optional) general funding statement \dots}%optional, to capture a funding statement, which applies to all authors. Please enter author specific funding statements as fifth argument of the \author macro.

\acknowledgements{I want to thank \dots}%optional

%\nolinenumbers %uncomment to disable line numbering



%Editor-only macros:: begin (do not touch as author)%%%%%%%%%%%%%%%%%%%%%%%%%%%%%%%%%%
\EventEditors{John Q. Open and Joan R. Access}
\EventNoEds{2}
\EventLongTitle{42nd Conference on Very Important Topics (CVIT 2016)}
\EventShortTitle{CVIT 2016}
\EventAcronym{CVIT}
\EventYear{2016}
\EventDate{December 24--27, 2016}
\EventLocation{Little Whinging, United Kingdom}
\EventLogo{}
\SeriesVolume{42}
\ArticleNo{23}
%%%%%%%%%%%%%%%%%%%%%%%%%%%%%%%%%%%%%%%%%%%%%%%%%%%%%%

%\usepackage[margin=0.7in]{geometry}
%\usepackage[margin=1in]{geometry}
%\usepackage{charter}

\usepackage{amsfonts, amsmath, amssymb, amsthm}
\usepackage{bm}
\usepackage{verbatim}
% \usepackage{color}
%\usepackage{euscript}
\usepackage{graphicx}
\usepackage[usenames,dvipsnames]{xcolor}
% BEGIN EDIT: Item 6 (underlining support)
% NEW:
\usepackage[normalem]{ulem}
% END EDIT: Item 6 (underlining support)
\newcommand{\deltext}[1]{{\color{red}\sout{#1}}\ }
\newcommand{\instext}[1]{{\color{ForestGreen}#1}}
\newenvironment{deltextblock}{\par\def\normalcolor{\color{red}}\color{red}\noindent\textbf{[Deleted:]}\ }{\par}
\newenvironment{instextblock}{\par\def\normalcolor{\color{ForestGreen}}\color{ForestGreen}\noindent\textbf{[Inserted:]}\ }{\par}
%\usepackage{paralist}
%\usepackage{subcaption}
\usepackage{tikz}
\usetikzlibrary{decorations.pathreplacing}
\usetikzlibrary{positioning}
\usetikzlibrary{arrows.meta}
\usetikzlibrary{decorations.pathmorphing}
\usetikzlibrary{calc}

\usepackage{algpseudocode}

\usepackage{xspace}
\usepackage{wrapfig}
\usepackage{url}




\def\figcapup{\vspace{-0mm}}
\def\figcapdown{\vspace{-0mm}}
\def\extraspacing{\vspace{2mm} \noindent}
\def\minilen{0.95\linewidth}
\def\vgap{\vspace{0.1 in}}



\newcommand{\apex}{\operatorname{apex}}

\newcommand{\eps}{\varepsilon}

\newcommand{\cF}{\mathcal{F}}
\newcommand{\cD}{\mathcal{D}}
\newcommand{\cO}{\mathcal{O}}
\newcommand{\cC}{\mathcal{C}}
\newcommand{\cG}{\mathcal{G}}
\newcommand{\cR}{\mathcal{R}}
\newcommand{\Ov}{\cO}


\newcommand{\tcr}[1]{\textcolor{red}{#1}}
\newcommand{\tcb}[1]{\textcolor{blue}{#1}}
\newcommand{\tcg}[1]{\textcolor{green}{#1}}
\newcommand{\tcc}[1]{\textcolor{cyan}{#1}}
\newcommand{\yakov}[1]{\textbf{Yakov: }\tcr{#1}}






% \author{}

\newcommand{\commented}[1]{}
% Old text (red, struck out) followed by the new text (green)
%\newcommand{\oldtext}[1]{{\color{red}\sout{#1}}\ }
\newcommand{\oldtext}{}
%\newcommand{\newtext}[1]{{\color{ForestGreen}#1}}
\newcommand{\newtext}[1]{#1}

\newenvironment{oldtextblock}{\par\def\normalcolor{\color{red}}\color{red}\noindent\textbf{[Old text:]}\ }{\par}
\newenvironment{newtextblock}{\par\def\normalcolor{\color{ForestGreen}}\color{ForestGreen}\noindent\textbf{[New text:]}\ }{\par}


\renewcommand{\arraystretch}{1.3}

% \newtheorem{theorem}{Theorem}
% \newtheorem{lemma}{Lemma}
% \newtheorem{corollary}{Corollary}
% \newtheorem{observation}{Observation}
% \newtheorem{proposition}{Proposition}
% \newtheorem{definition}{Definition}
% \newtheorem{problem}{Problem}

\begin{document}
\maketitle
\begin{abstract}
In this paper we describe new data structures for  adaptive and approximate orthogonal range counting in three dimensions.   
\end{abstract}
\section{Preliminaries}
\label{sec:prelim}
For a three-dimensional geometric object $x$ (point, rectangle, etc.) we denote by $\pi(x)$ the projection of $x$ onto the $xy$-plane. 

shallow cuttings, definition

We associate a polygon $R(C)$ with every cell $C$ of $\cC$ in the following way. We visit  the apex points of cells $C\in \cC$ in the decreasing order of their $z$-coordinates; for every point $r=\mathrm{apex}(C)$, we shoot vertical rays from $\pi(r)$ until they hit the axes or one of the staircases of $\cC$, or one of the previously constructed polygons. We will denote the polygon of $C$ by $R(C)$.   

\section{Dominance Queries}
\label{sec:domin-approx}
\subsection{Step 1}
\label{sec:step1}
Let $t_i=n^{g(i)}$ where $g(i)=(\frac{2}{3})^i$. We construct a sequence of $t_i$-shallow cuttings $\cC_i$ and  planar  subdivisions $\cR_i$ for $i=0$, $1$,$\ldots$.

Since $t_0=n$, the $t_0$-shallow cutting consists of a single cell $C$  that contains all points of the input set $P$, so that $L(C)=P$. The subdivision $\cR_0$ contains a single rectangle $R_0=\pi(C)$.  The $t_1$-shallow cutting consists of $n^{1/3}$ cells that cover the $\leq n^{2/3}$-level of $P$. We associate a polygon $R_j=R(C_j)$ to every cell $C_j$ of $\cC_1$. 

Suppose that the $t_{i-1}$-shallow cutting $\cC_i$ is already constructed. We construct $\cC_i$ and $\cR_i$ as follows. For each cell $C_j\in \cC_{i-1}$, we construct an $n^{g(i)}$-shallow cutting  $\cG^i_j$.  A cell $C$ of $\cG^i_j$ is \emph{valid} iff $\pi(\apex(C))$ is in $R_j$.
We visit all valid cells $C_k$ of $\cG^i_j$ in the decreasing order of their $z$ coordinates and let $R_k=R(C_k)\cap R_j$. The $t_i$-shallow cutting $\cC_i$ is the union of all  valid cells $C_j$ in all $\cG^i_j$. The set $\cR_i$ consists of all polygons $R_k$ for all cells $C_k$ in $\cC_i$. 

\begin{lemma}
  $\cC_i$ has  $O(\frac{n}{n^{g(i)}}c^i$ cells for some constant $c$.  $\cC_i$ covers the $\leq t_i$-level of $P$.
\end{lemma}
\begin{proof}
  Let $r$ denote arbitrary point, such that the level of $r$ is at most $t_i$.  Since the $\leq t_i$-level of $P$ is contained in the $t_{i-1}$-level, $r$ is in at least one cell of $\cC_{i-1}$. Let $C_k$ denote the cell of $\cC_{i-1}$ with the highest $z$-coordinate such that $\apex(C_k)$ dominates $r$.  Then $\pi(r)$ is in the polygon of $C_k$.  Let $G_k$ denote the cell of $\cG^i_k$ with the highest $z$-coordinate, such that $G_k$ contains $r$. Since   $\apex(G_k)$ dominates $r$  and $\apex(C_k)$ dominates $\apex(G_k)$,  $\pi(\apex(G_k))$  is in the polygon of $C_k$. Hence $G_k$ is a valid cell of $\cG^i_k$ and $G_k$ is in $\cC_i$.
\end{proof}

\begin{lemma}
  \label{lemma:subdiv-size}
  $\cR_i$ is a subdivision of the plane with $O(|\cC_i|)$ vertices. 
\end{lemma}
\begin{proof}
  
\end{proof}


\begin{lemma}
  \label{lemma:rect-contain}
  Each polygon $R^i\in \cR_i$ is contained in some polygon $R^{i-1}\in \cR_{i-1}$.  
\end{lemma}
\begin{proof}
  Directly follows from the construction: every polygon of a cell $G\in G^i_k$ is a sub-polygon of $R_k$, where $R_k$ is the polygon of $C_k\in \cC_{i-1}$.
\end{proof}

We construct shallow cuttings $\cC_i$ and planar subdivisions $\cR_i$ for $i=1,2,\ldots g$, such that $n^{g(i)}\leq \log^{c'}n$ for some constant $c'$ that will be defined later.  For each polygon $R\in \cR_g$, we store the polygons $R^j\in \cR_j$, $0\le j < g$, such that $R^j$ contains $R$;
for every polygon $R^j$ we also store the corresponding cell $C^j\in \cC_j$.  For a given query point $q$, we identify the rectangle $R_g\in \cR_g$ that contains $\pi(q)$.   Then we examine all cells $C^j$ and find the largest $j$ such that $C^j$ contains $q$.  Since we can locate $\pi_q$ in $\cR_g$ in $O(\log\log n)$ time using the data structure of Chan~\cite{Chan??}, the overall query time is $O(\log \log n)$.

\subsection{Step 2: Better Approximation}
\label{sec:step2}
For each cell $C$ of a $t_i$-shallow cutting $\cC_i$ we construct a sequence of $f_j$-shallow cuttings $\cF_j(C)$ for $L(C)$. We set  $f_j=t_i/(c_1)^j$ for a constant $c_1$ that will be defined later.  We consider the planar subdivision $\cR_j(C)$ for each $\cF_j(C)$ and the overlay $Ov(C)$ of all $\cR_j(C)$.  has complexity $O(t_i^{2/3}\log t_i)=o(t_i)$: since $F_j(C)$ has $O(t_i^{1/3})$ cells, $\cR_j(C)$ has $O(t_i^{1/3})$ segments. Since a segment of $R_j(C)$ can intersect $O(t_i^{1/3})$ other segments, the complexity of $\Ov(C)$ is $O(t_i^{2/3})$. We store the point location data structure for $\Ov(C)$.  Additionally, for every face $g$  of the overlay $\Ov(C)$, we store the list of cells $I(g)$; $I(g)$ contains cells $F_j\in \cF_j(C)$, such that the polygon of $F_j$ contains $g$. Cells in $I(g)$ are ordered in decreasing order of $j$. 


Suppose that the query point $q$ is contained in a cell $C$ of the $t_i$-shallow cutting.  And we know that $q$ is not contained in any cell of a $t_{i+1}$-shallow cutting, i.e., the level of $q$ is larger than $t_{i+1}$. We find the face $g_q$ of $\Ov(C)$ in $O(\log\log n)$ time. Then we use binary search on $I(g_q)$ and find the largest $j$, such that $F_j(C)$ contains $q$.  The overall query time is $O(\log\log n)$.

When $j$ is found, we obtain a constant-factor approximation for the number of points dominated by $q$: $q$ dominates more than $\frac{t_i}{c_1^{j+1}}$ points and less than $\frac{t_i}{c_1^j}$ points.  

\subsection{Step 3: Refining}
\label{sec:step3}
We keep a sequence of $(\delta)^i$-shallow cuttings $\cD_i$. We calculate the largest $k$, such that $(\delta)^k\leq \frac{t_i}{c_1^{j+1}}$. We increment $k$ by $1$ until we find the smallest index $k$, such that $q$ is in some cell of $\cD_k$. The index $k$ is incremented $\frac{c_1}{\delta}=O(1)$ times. For every $\cD_k$, we can check if $q$ is in some cell of $\cD_k$ in $O(\log\log n)$ time. Hence the overall time of the third step is $O(\log \log n)$.

Our result can be summed up in the following theorem.
\begin{theorem}
  \label{theor:domin-approx}
 There exists a data structure that answers approximate dominance reporting queries in $O(\log \log n)$ time. The data structure uses $O(n)$ space.
\end{theorem}


\section{Approximate Range Counting for General 3-d Queries}
\label{sec:label}
\begin{lemma}
  There exists an $O(n\log^{\eps}n)$-space data structure that answers 5-sided 3-d approximate range counting queries in $O(\log \log n)$ time. 
\end{lemma}

\begin{theorem}
  \label{theor:3d-approx}
  There exists an $O(n\log^{1+\eps}n)$-space data structure that answers general 3-d approximate range counting queries in $O(\log \log n)$ time. 
\end{theorem}




\end{document}


%%% Local Variables:
%%% mode: LaTeX
%%% TeX-master: t
%%% End:
